\documentclass[12pt,a4paper]{article}
\usepackage[utf8]{inputenc}
\usepackage[T1]{fontenc}
\usepackage{lmodern}
\usepackage{amsmath,amssymb,amsthm}
\usepackage{amsfonts}
\usepackage{mathrsfs}
\usepackage{geometry}
\usepackage{titlesec}
\usepackage{fancyhdr}
\usepackage{array}
\usepackage{longtable}

\geometry{margin=1in}

\pagestyle{fancy}
\fancyhf{}
\fancyhead[C]{\thepage}
\renewcommand{\headrulewidth}{0pt}

\titleformat{\section}{\large\bfseries}{\thesection.}{0.5em}{}
\titleformat{\subsection}{\normalsize\bfseries}{\thesubsection.}{0.5em}{}

% Calligraphic / fraktur-style letters used by Cayley for his special functions
\newcommand{\fS}{\mathfrak{S}}
\newcommand{\fA}{\mathfrak{A}}
\newcommand{\fB}{\mathfrak{B}}
\newcommand{\fC}{\mathfrak{C}}
\newcommand{\fD}{\mathfrak{D}}
\newcommand{\fG}{\mathfrak{G}}
% Helper notation for the index pairs (m,n), (\bar m,n), (m,\bar n), (\bar m, \bar n)
\newcommand{\mn}{(m,n)}
\newcommand{\mbn}{(\bar m,n)}
\newcommand{\mnb}{(m,\bar n)}
\newcommand{\mbnb}{(\bar m,\bar n)}

\usepackage{graphicx}
\emergencystretch=3em
\tolerance=600
\begin{document}

\begin{center}
{\LARGE\bfseries The Collected Mathematical Papers of}\\[0.5em]
{\LARGE\bfseries Arthur Cayley}\\[1em]
{\Large Volume I}\\[1em]
{\large Pages 126--150}\\[2em]
\end{center}

\tableofcontents
\newpage

% ===========================================================================
% Page 126 -- continuation of paper [20] On Certain Results Relating to Quaternions
% ===========================================================================
\section*{[20] (concluded)}
\section*{On Certain Results Relating to Quaternions.}
\addcontentsline{toc}{section}{20. On Certain Results Relating to Quaternions (conclusion)}

\textit{[From the Philosophical Magazine, vol.~xxvi.~(1845), pp.~141--145.]}

\bigskip
\setcounter{page}{126}

On the contrary, the result of this elimination is the very different equation
\begin{equation*}
M\psi^{-1}\cdot\varpi - M\varpi^{-1}\cdot\psi = 0 \quad \cdots(15),
\end{equation*}
equivalent of course to four independent equations, one of which may evidently be replaced by
\begin{equation*}
M\psi\cdot M\varpi' - M\varpi\cdot M\psi' = 0 \quad \cdots(16),
\end{equation*}
if $M\varpi$, \&c.\ denotes the modulus of $\varpi$, \&c. An equation analogous to this last will undoubtedly hold for any number of equations, but it is difficult to say what is the equation analogous to the one immediately preceding this, in the case of a greater number of equations, or rather, it is difficult to give the result in a symmetrical form independent of extraneous factors.

I may just, in conclusion, mention what appears to me a possible application of Sir William Hamilton's interesting discovery. In the same way that the circular functions depend on infinite products, such as
\begin{equation*}
\sin x = x\prod\left(1 + \frac{x}{m\pi}\right),\quad\&\text{c.} \cdots(17),
\end{equation*}
($m$ any integer from $\infty$ to $-\infty$, omitting $m=0$)
and the inverse elliptic functions on the doubly infinite products
\begin{equation*}
x\prod\prod\left(1 + \frac{x}{m\varpi + n\varpi'i}\right) \cdots(18),
\end{equation*}
($m$ and $n$ integers from $\infty$ to $-\infty$, omitting $m=0$, $n=0$),
may not the inverse ultra-elliptic functions of the next order of complexity depend on the quadruply infinite products
\begin{equation*}
x\prod\prod\prod\prod\left(1 + \frac{x}{m\varpi + n\varpi' i + o\rho j + p\rho' k}\right) \cdots(19),
\end{equation*}
($m$, $n$, $o$, $p$ integers from $\infty$ to $-\infty$, omitting $m=0$, $n=0$, $o=0$, $p=0$).

It seems as if some supposition of this kind would remove a difficulty started by Jacobi (Crelle, t.~ix.) with respect to the multiple periodicity of these functions. Of course this must remain a mere suggestion until the theory of quaternions is very much more developed than it is at present; in particular the theory of quaternion exponentials would have to be developed, for even in a product, such as (18), there is a certain singular exponential factor running through the theory, as appears from some formulae in Jacobi's \textit{Fund.\ Nova} (relative to his functions $\Theta$, $H$), the occurrence of which may be accounted for, \textit{\`a priori}, as I have succeeded in doing in a paper to be published shortly in the Cambridge Mathematical Journal [24].

% ===========================================================================
% Page 127 -- Paper [21]
% ===========================================================================
\newpage
\section*{[21]}
\section*{On Jacobi's Elliptic Functions, in reply to the Rev.~B.~Bronwin; and on Quaternions.}
\addcontentsline{toc}{section}{21. On Jacobi's Elliptic Functions, in reply to the Rev.~B.~Bronwin; and on Quaternions}

\textit{[From the Philosophical Magazine, vol.~xxvi.~(1845), pp.~208, 211.]}

\bigskip
\setcounter{page}{127}

The first part of this Paper is omitted, see [17]: only the Postscript on Quaternions, pp.~210, 211, is printed.

\medskip

It is possible to form an analogous theory with seven imaginary roots of $(-1)$ (with $v = 2^n - 1$ roots when $v$ is a prime number). Thus if these be $i_1, i_2, i_3, i_4, i_5, i_6, i_7$, which group together according to the types
\[
123,\ 145,\ 624,\ 653,\ 725,\ 734,\ 176,
\]
i.e.\ the type $123$ denotes the system of equations
\begin{align*}
i_1 i_2 &= i_3, & i_2 i_3 &= i_1, & i_3 i_1 &= i_2,\\
i_2 i_1 &= -i_3, & i_3 i_2 &= -i_1, & i_1 i_3 &= -i_2,
\end{align*}
\&c.\ We have the following expression for the product of two factors:
\begin{align*}
&(x_0 + x_1 i_1 + \ldots + x_7 i_7) \times (x'_0 + x'_1 i_1 + \ldots + x'_7 i_7)\\
&\quad= (01) + [12 + 45 + 76 + (01)]\,i_1 + [13 + 46 + 57 + (02)]\,i_2\\
&\quad+ [14 + 65 + 76 + (03)]\,i_3 + [51 + 36 + 47 + (04)]\,i_4\\
&\quad+ [14 + 53 + 72 + (05)]\,i_5 + [24 + 35 + 71 + (06)]\,i_6\\
&\quad+ [25 + 34 + 61 + (07)]\,i_7,
\end{align*}
where $(01) = x_0 x'_0 + x_1 x'_1 + \ldots$;\ \ $12 = x_1 x'_2 - x_2 x'_1$, \&c.;
and the modulus of this expression is the product of the moduli of the factors. The above system of types requires some care in writing down, and not only with respect to the combinations of the letters, but also to their order; it would be vitiated, e.g.\ by writing $716$ instead of $176$. A theorem analogous to that which I gave before, for quaternions, is the following:---If $A = 1 + \lambda_1 i_1 \ldots + \lambda_7 i_7$, $X = x_1 i_1 \ldots + x_7 i_7$: it is immediately shown that the possible part of $A^{-1}XA$ vanishes, and that the coefficients of $i_1, \ldots, i_7$ are linear functions of $x_1, \ldots, x_7$. The modulus of the above expression is evidently the modulus of $X$; hence ``we may determine seven linear functions of $x_1\ldots x_7$, the sum of whose squares is equal to $x_1^2 + \ldots + x_7^2$.'' The number of arbitrary quantities is however only seven, instead of twenty-one, as it should be.

% ===========================================================================
% Page 128 -- Paper [22]
% ===========================================================================
\newpage
\section*{[22]}
\section*{On Algebraical Couples.}
\addcontentsline{toc}{section}{22. On Algebraical Couples}

\textit{[From the Philosophical Magazine, vol.~xxvii.~(1845), pp.~38--40.]}

\bigskip
\setcounter{page}{128}

It is worth while, in connection with the theory of quaternions and the researches of Mr Graves (Phil.\ Mag.\ [Vol.~xxvi.~(1845), pp.~315--320]), to investigate the properties of a couple $ix + jy$ in which $i$, $j$ are symbols such that
\begin{align*}
i^2 &= \alpha i + \beta j,\\
ij &= \alpha' i + \beta' j,\\
ji &= \gamma i + \delta j,\\
j^2 &= \gamma' i + \delta' j.
\end{align*}
If $(ix + jy)(ix_1 + jy_1) = iX + jY$, then
\begin{align*}
X &= \alpha xx_1 + \alpha' xy_1 + \gamma x_1 y + \gamma' y y_1,\\
Y &= \beta xx_1 + \beta' xy_1 + \delta x_1 y + \delta' y y_1.
\end{align*}
Imagine the constants $\alpha$, $\beta$, $\ldots$ so determined that $ix + jy$ may have a modulus of the form $K(x + \lambda y)(x + \mu y)$; there results one of the four following essentially independent systems:

\medskip
\textbf{A.}
\begin{align*}
i^2 &= -\tfrac{1}{\lambda\mu}\delta i - \beta j,\\
ij &= \alpha' i + \beta' j,\\
ji &= \gamma i + \delta j,\\
j^2 &= -\lambda\mu\beta i + (\gamma + \lambda\delta + \mu\delta')j,\\
X + \lambda Y &= -(\gamma + \lambda\delta)(x + \lambda y)(x_1 + \lambda y_1),\\
X + \mu Y &= -(\gamma + \mu\delta)(x + \mu y)(x_1 + \mu y_1).
\end{align*}

% ===========================================================================
% Page 129
% ===========================================================================
\newpage
\setcounter{page}{129}

The couple may be said to have the two linear moduli,
\[
X + \lambda Y,\quad X + \mu Y,
\]
as well as the quadratic one,
\[
-(\gamma + \lambda\delta)(\gamma + \mu\delta)(x + \lambda y)(x + \mu y),
\]
the product of these, which is the modulus, and the only modulus in the remaining systems.

\medskip
\textbf{B.}
\begin{align*}
i^2 &= -\beta i - \tfrac{\gamma}{\mu}j,\\
ij &= ji = \gamma i + \delta j,\\
j^2 &= (\lambda + \mu)\gamma i - \lambda\mu\delta\, j,\\
X + \lambda Y &= \tfrac{1}{\mu-\lambda}(\gamma + \lambda\delta)(x + \mu y)(x_1 + \mu y_1),\\
X + \mu Y &= \tfrac{1}{\mu-\lambda}(\gamma + \mu\delta)(x + \mu y)(x_1 + \mu y_1).
\end{align*}

\medskip
\textbf{C.}
\begin{align*}
i^2 &= -\tfrac{\mu^2\beta + \lambda^2}{\mu - \lambda}\,j + \cdots,\\
ji &= \gamma i + \delta j,\\
j^2 &= (\lambda + \mu)\gamma i + \lambda\mu\delta\, j,\\
X + \lambda Y &= \tfrac{1}{\mu-\lambda}(\gamma + \lambda\delta)(x + \lambda y)(x_1 + \mu y_1),\\
X + \mu Y &= -\tfrac{1}{\mu-\lambda}(\gamma + \mu\delta)(x + \mu y)(x_1 + \lambda y_1).
\end{align*}

\medskip
\textbf{D.}
\begin{align*}
i^2 &= -\beta i + \cdots,\\
ji &= \gamma i + \delta j,\\
j^2 &= -\lambda\mu\beta i + (\gamma + \lambda + \mu\delta)j,\\
X + \lambda Y &= -(\gamma + \lambda\delta)(x + \mu y)(x_1 + \lambda y_1),\\
X + \mu Y &= -(\gamma + \mu\delta)(x + \lambda y)(x_1 + \mu y_1).
\end{align*}

% ===========================================================================
% Page 130
% ===========================================================================
\newpage
\setcounter{page}{130}

The formulae are much simpler and not essentially less general, if $\mu = -\lambda$. They thus become

\medskip
\textbf{A$'$.}
\begin{align*}
i^2 &= \delta i + \beta j,\\
ij &= ji = \gamma i + \delta j,\\
j^2 &= \lambda^2\beta i,\\
X \pm \lambda Y &= \pm(\gamma \pm \lambda\delta)(x \pm \lambda y)(x_1 \pm \lambda y_1).
\end{align*}
(Two linear moduli.)

\medskip
\textbf{B$'$.}
\begin{align*}
i^2 &= -\beta i - \tfrac{\gamma}{\lambda}j,\\
ij &= ji = \gamma i + \delta j,\\
j^2 &= -\lambda^2\beta i - \gamma j,\\
X \pm \lambda Y &= \mp\tfrac{1}{2\lambda}(\gamma \pm \lambda\delta)(x \pm \lambda y)(x_1 \mp \lambda y_1).
\end{align*}

\medskip
\textbf{C$'$.}
\begin{align*}
i^2 &= -\gamma i - \beta j,\\
ji &= \gamma i + \delta j,\\
j^2 &= -\lambda^2\delta i - \gamma j,\\
X \pm \lambda Y &= \pm\tfrac{1}{2\lambda}(\gamma \pm \lambda\delta)(x \pm \lambda y)(x_1 \mp \lambda y_1).
\end{align*}

\medskip
\textbf{D$'$.}
\begin{align*}
i^2 &= -\delta i - \beta j,\\
ij &= -\gamma i - \delta j,\\
ji &= \gamma i + \delta j,\\
j^2 &= \lambda^2\beta i + \gamma j,\\
X \pm \lambda Y &= \mp\tfrac{1}{2\lambda}(\gamma + \lambda\delta)(x + \lambda y)(x_1 \mp \lambda y).
\end{align*}

There is a system more general than (A.) having a single linear modulus $q(\theta x + Y)$: this is

\medskip
\textbf{E.}
\begin{align*}
i^2 &= \alpha(i - \theta j) + \theta^2 q j,\\
ij &= \alpha'(i - \theta j) + \theta q j,\\
ji &= \gamma(i - \theta j) + \theta q j,\\
j^2 &= \gamma'(i - \theta j) + q j,\\
\theta X + Y &= q(\theta x + y)(\theta x_1 + y_1);
\end{align*}
or, without real loss of generality,

% ===========================================================================
% Page 131
% ===========================================================================
\newpage
\setcounter{page}{131}

\medskip
\textbf{E$'$.}
\begin{align*}
i^2 &= \alpha i,\\
ij &= \alpha' i,\\
ji &= \gamma i,\\
j^2 &= \gamma i + q j,\\
Y &= q y y_1.
\end{align*}

To complete the theory of this system, one may add the identical equation
\[
\frac{1}{\alpha - q\omega}\cdot\frac{1}{\alpha - q\omega_1}
= -\theta^2 \frac{1}{(\alpha' + \alpha'\omega)(\alpha' + \alpha'\omega_1)} \cdot (\alpha + \alpha\omega)(\alpha + \alpha\omega_1),
\]
where
\[
\omega + \omega_1 = -\frac{\alpha' - \theta\gamma}{\alpha - q\theta},\quad
\omega\omega_1 = \frac{\alpha' - \theta\gamma'}{\alpha - q\theta}.
\]
By determining the constants, so that
\[
\frac{\alpha'}{\alpha} = \frac{\alpha' - \theta\gamma}{\alpha - q\theta} = \frac{\alpha' - \theta\gamma'}{\alpha' - q\theta}\,\omega_1,
\]
the system would reduce itself to the form A.

% ===========================================================================
% Page 132 -- Paper [23]
% ===========================================================================
\newpage
\section*{[23]}
\section*{On the Transformation of Elliptic Functions.}
\addcontentsline{toc}{section}{23. On the Transformation of Elliptic Functions}

\textit{[From the Philosophical Magazine and Journal of Science, vol.~xxvii.~(1845), pp.~424--427.]}

\bigskip
\setcounter{page}{132}

In a former paper [19] I gave a proof of Jacobi's theorem, which I suggested [21] would lead to the resolution of the very important problem of finding the relation between the complete functions. This is in fact effected by the formulae there given, but there is an apparent indeterminateness in them, the cause of which it is necessary to explain, and which I shall now show to be inherent in the problem. For the sake of supplying an omission, for the detection of which I am indebted to Mr Bronwin, I will first recapitulate the steps of the demonstration.

If $\tfrac{1}{2}\omega$, $\tfrac{1}{2}\nu$\footnote{Analogous to the $K$, $K'i$ of M.~Jacobi.} be the complete functions corresponding to $\phi$, $x$, then this function is expressible in the form
\[
\phi x = x \prod\prod \frac{1 + \dfrac{x}{m\omega + n\nu}}{1 + \dfrac{x}{m + \tfrac{1}{2}\omega + n + \tfrac{1}{2}\nu}}\,.
\]
Let $p$ be any prime number, $\mu$, $\nu$ integers not divisible by $p$, and
\[
\theta = \frac{\mu\omega + \nu\nu}{p}.
\]
The function
\[
\phi_1 x = \frac{\phi(x+2\theta)\,\phi(x+4\theta)\,\ldots\,\phi(x+2(p-1)\theta)}{\phi(2\theta)\,\ldots\,\phi(2(p-1)\theta)}
\]
is always reducible to the form
\[
x\prod\prod \frac{1 + \dfrac{x}{m'\omega' + n'\nu'}}{1 + \dfrac{x}{m'+\tfrac{1}{2}\omega' + n'+\tfrac{1}{2}\nu'}}\,,
\]

% ===========================================================================
% Page 133
% ===========================================================================
\newpage
\setcounter{page}{133}

and thus $\phi_1 x$ is an inverse function, the complete functions of which are $\tfrac{1}{2}\omega'$, $\tfrac{1}{2}\nu'$: and $\omega'$, $\nu'$ are connected with $\omega$, $\nu$ by the equations
\begin{align*}
\omega' &= \tfrac{1}{p}(\alpha\omega + \beta\nu),\\
\nu' &= \tfrac{1}{p}(\alpha'\omega + \beta'\nu),
\end{align*}
$\alpha$, $\beta$, $\alpha'$, $\beta'$ being any integers subject to the conditions that $\alpha$, $\beta'$ are odd and $\alpha'$, $\beta$ even; also
\begin{align*}
\alpha\beta' - \alpha'\beta &= p,\\
\mu\beta' - \nu\alpha' &= l'p,\\
\mu\beta - \nu\alpha &= lp,
\end{align*}
$l$, $l'$ being any integers whatever. In fact, to prove this, we have only to consider the general form of a factor in the numerator of $\phi_1 x$. Omitting a constant factor, this is
\[
\left(1 + \frac{x}{m\omega + n\nu + 2r\theta}\right),
\]
and it is to be shown that we can always satisfy the equation
\[
m\omega + n\nu + 2r\theta = m'\omega' + n'\nu',
\]
or the equations
\begin{align*}
pm + 2r\mu &= m'\alpha + n'\alpha',\\
pn + 2r\nu &= m'\beta + n'\beta';
\end{align*}
and also that to each set of values of $m$, $n$, $r$, there is a unique set of values of $m'$, $n'$, and vice versa. This is done in the paper referred to. Moreover, with the suppositions just made as to the numbers $\alpha$, $\beta'$ being odd and $\alpha'$, $\beta$ even, it is obvious that $m'$ is odd or even, according as $m$ is, and $n'$ according as $n$ is, which shows that we can likewise satisfy
\[
m + \tfrac{1}{2}\omega + n + \tfrac{1}{2}\nu + 2r\theta = m' + \tfrac{1}{2}\omega' + n' + \tfrac{1}{2}\nu';
\]
and thus the denominator of $\phi_1 x$ is also reducible to the required form.

Now proceeding to the immediate object of this paper, $\alpha$, $\beta$, $\alpha'$, $\beta'$, and consequently $\omega'$, $\nu'$, are to a certain extent indeterminate. Let $A$, $B$, $A'$, $B'$ be a particular set of values of $\alpha$, $\beta$, $\alpha'$, $\beta'$, and $\Omega$, $\Upsilon$ the corresponding values of $\omega'$, $\nu'$. We have evidently $A$, $B'$ odd and $A'$, $B$ even. Also
\begin{align*}
AB' - A'B &= p,\\
\mu B' - \nu A' &= L'p,\\
\mu B - \nu A &= Lp,\\
\Omega &= \tfrac{1}{p}(A\omega + B\nu),\\
\Upsilon &= \tfrac{1}{p}(A'\omega + B'\nu).
\end{align*}

% ===========================================================================
% Page 134
% ===========================================================================
\newpage
\setcounter{page}{134}

By eliminating $\omega$, $\nu$ from these equations and the former system, it is easy to obtain
\begin{align*}
\omega' &= a\Omega + b\Upsilon,\\
\nu' &= a'\Omega + b'\Upsilon,
\end{align*}
where
\begin{align*}
a &= \tfrac{1}{p}(\alpha B' - \beta A'), & b &= -\tfrac{1}{p}(\alpha B - \beta A),\\
a' &= \tfrac{1}{p}(\alpha' B' - \beta' A'), & b' &= -\tfrac{1}{p}(\alpha' B - \beta' A).
\end{align*}
The coefficients $a$, $b$, $a'$, $b'$ are integers, as is obvious from the equation
\[
\mu(\alpha B - \beta A) = p(L'\alpha - lA'),
\]
and the others analogous to it; moreover, $a$, $b'$ are odd and $a'$, $b$ even, and
\[
ab' - a'b = \tfrac{1}{p^2}(AB' - A'B)(\alpha\beta' - \alpha'\beta);
\]
that is
\[
ab' - a'b = 1.
\]

Hence the theorem,---``The general values $\omega'$, $\nu'$ of the complete functions are linearly connected with the particular system of values $\Omega$, $\Upsilon$ by the equations, $\omega' = a\Omega + b\Upsilon$, $\nu' = a'\Omega + b'\Upsilon$, in which $a$, $b'$ are odd integers and $a'$, $b$ even ones, satisfying the condition $ab' - a'b = 1$.''

With this relation between $\Omega$, $\Upsilon$ and $\omega'$, $\nu'$, it is easy to show that the function $\phi_1 x$ is precisely the same, whether $\Omega$, $\Upsilon$ or $\omega'$, $\nu'$ be taken for the complete functions. In fact, stating the proposition relatively to $\phi x$, we have,---``The inverse function $\phi x$ is not altered by the change of $\omega$, $\nu$ into $\omega'$, $\nu'$, where $\omega' = a\omega + b\nu$, $\nu' = a'\omega + b'\nu$, and $a$, $b$, $a'$, $b'$ satisfy the conditions that $a$, $b'$ are odd, $a'$, $b$ even, and $ab' - a'b = 1$.'' This is immediately shown by writing
\[
m\omega + n\nu = m'\omega' + n'\nu',
\]
or
\begin{align*}
m &= m'a + n'a',\\
n &= m'b + n'b'.
\end{align*}
It is obvious that to each set of values of $m$, $n$ there is a unique set of values of $m'$, $n'$, and vice vers\^a: also that odd or even values of $m$, $m'$ or $n$, $n'$ always correspond to each. It is, in fact, the preceding reasoning applied to the case of $p = 1$.

Hence finally the theorem,---``The only conditions for determining $\omega'$, $\nu'$ are the equations
\[
\omega' = \tfrac{1}{p}(\alpha\omega + \beta\nu),\quad \nu' = \tfrac{1}{p}(\alpha'\omega + \beta'\nu),
\]
where $\alpha$, $\beta'$ are odd and $\alpha'$, $\beta$ even, and
\[
\alpha\beta' - \alpha'\beta = p,\quad \mu\beta' - \nu\alpha' = l'p,\quad \mu\beta - \nu\alpha = lp,
\]
$l$ and $l'$ arbitrary integers: and it is absolutely indifferent what system of values is adopted for $\omega'$, $\nu'$, the value of $\phi_1 x$ is precisely the same.''

% ===========================================================================
% Page 135
% ===========================================================================
\newpage
\setcounter{page}{135}

We derive from the above the somewhat singular conclusion, that the complete functions are not absolutely determinate functions of the modulus; notwithstanding that they are given by the apparently determinate conditions,
\[
\tfrac{1}{2}\omega = \int_0^{\frac{1}{c}} \frac{dx}{\sqrt{(1 - c^2 x^2)(1 + e^2 x^2)}},\quad
\tfrac{1}{2}\nu = \int_0^{\frac{1}{e}} \frac{dx}{\sqrt{(1 + c^2 x^2)(1 - e^2 x^2)}}.
\]

In fact definite integrals are in many cases really indeterminate, and acquire different values according as we consider the variable to pass through real values, or through imaginary ones. Where the limits are real, it is tacitly supposed that the variable passes through a succession of real values, and thus $\omega$, $\nu$ may be considered as completely determined by these equations, but only in consequence of this tacit supposition. If $c$ and $e$ are imaginary, there is absolutely no system of values to be selected for $\omega$, $\nu$ in preference to any other system. The only remaining difficulty is to show from the integral itself, independently of the theory of elliptic functions, that such integrals contain an indeterminateness of two arbitrary integers; and this difficulty is equally great in the simplest cases. Why, \textit{\`a priori}, do the functions
\[
\sin^{-1} x = \int_0^x \frac{dx}{\sqrt{1-x^2}},\quad \text{or}\quad \log x = \int \frac{dx}{x}
\]
contain a single indeterminate integer?

I am of course aware, that in treating of the properties of such products as $\prod\!\!\prod\!\left(1 + \dfrac{x}{m\omega + n\nu}\right)$, it is absolutely necessary to pay attention to the relations between the infinite limiting values of $m$ and $n$; and that this introduces certain exponential factors, to which no allusion has been made. But these factors always disappear from the quotient of two such products, and to have made mention of them would only have been embarrassing the demonstration without necessity.

% ===========================================================================
% Page 136 -- Paper [24]
% ===========================================================================
\newpage
\section*{[24]}
\section*{On the Inverse Elliptic Functions.}
\addcontentsline{toc}{section}{24. On the Inverse Elliptic Functions}

\textit{[From the Cambridge Mathematical Journal, t.~iv.~(1845), pp.~257--277.]}

\bigskip
\setcounter{page}{136}

The properties of the inverse elliptic functions have been the object of the researches of the two illustrious analysts, Abel and Jacobi. Among their most remarkable ones may be reckoned the formulae given by Abel (\OE uvres, t.~i.\ p.~212 [Ed.~2, p.~343]), in which the functions $\phi a$, $fa$, $Fa$, (corresponding to Jacobi's $\sin\operatorname{am} a$, $\cos\operatorname{am} a$, $\Delta\operatorname{am} a$, though not precisely equivalent to these, Abel's radical being $[(1-c^2x^2)(1+e^2x^2)]^{\frac{1}{2}}$, and Jacobi's, like that of Legendre's $[(1-x^2)(1-k^2 x^2)]^{\frac{1}{2}}$), are expressed in the form of fractions, having a common denominator; and this, together with the three numerators, resolved into a doubly infinite series of factors; i.e.\ the general factor contains two independent integers. These formulae may conveniently be referred to as ``Abel's double factorial expressions'' for the functions $\phi$, $f$, $F$. By dividing each of these products into an infinite number of partial products, and expressing these by means of circular or exponential functions, Abel has obtained (pp.~216--218) two other systems of formulae for the same quantities, which may be referred to as ``Abel's first and second single factorial systems.'' The theory of the functions forming the above numerators and denominator, is mentioned by Abel in a letter to Legendre (\OE uvres, t.~ii.\ p.~259 [Ed.~2, p.~272]), as a subject to which his attention had been directed, but none of his researches upon them have ever been published. Abel's double factorial expressions have nowhere anything analogous to them in Jacobi's \textit{Fund.\ Nova}; but the system of formulae analogous to the first single factorial system is given by Jacobi (p.~86), and the second system is implicitly contained in some of the subsequent formulae. The functions forming the numerator and denominator of $\sin\operatorname{am} u$, Jacobi represents, omitting a constant factor, by $H(u)$, $\Theta(u)$; and proceeds to investigate the properties of these new functions. This he principally effects by means of a very remarkable equation of the form
\[
\frac{1}{\Theta(u)}\cdot\frac{d\Theta(u)}{du} = \tfrac{1}{2}Au + B\int_0^u du \cdot \int_0^u du \sin^2\operatorname{am} u,
\]
(\textit{Fund.\ Nova}, pp.~145, 133), by which $\Theta(u)$ is made to depend on the known function $\sin\operatorname{am} u$. The other two numerators are easily expressed by means of the two functions $H$, $\Theta$.

% ===========================================================================
% Page 137
% ===========================================================================
\newpage
\setcounter{page}{137}

From the omission of Abel's double factorial expressions, which are the only ones which display clearly the real nature of the functions in the numerators and denominators; and besides, from the different form of Jacobi's radical, which complicates the transformation from an impossible to a possible argument, it is difficult to trace the connection between Jacobi's formulae; and in particular to account for the appearance of an exponential factor which runs through them. It would seem therefore natural to make the whole theory depend upon the definitions of the new transcendental functions to which Abel's double factorial expressions lead one, even if these definitions were not of such a nature, that one only wonders they should never have been assumed \textit{\`a priori} from the analogy of the circular functions $\sin$, $\cos$, and quite independently of the theory of elliptic integrals. This is accordingly what I have done in the present paper, in which therefore I assume no single property of elliptic functions, but demonstrate them all, from my fundamental equations. For the sake however of comparison, I retain entirely the notation of Abel. Several of the formulae that will be obtained are new.

\medskip
The infinite product
\[
x\prod\left(1 + \frac{x}{m\omega}\right) \cdots(1),
\]
where $m$ receives the integer values $\pm1$, $\pm2$,\ldots $\pm r$, converges, as is well known, as $r$ becomes indefinitely great to a determinate function $\sin\frac{\pi x}{\omega}$ of $x$; the theory of which might, if necessary, be investigated from this property assumed as a definition. We are thus naturally led to investigate the properties of the new transcendant
\[
u = x\prod\prod\left(1 + \frac{x}{m\omega + n\varpi i}\right) \cdots(2):
\]
$m$ and $n$ are integer numbers, positive or negative; and it is supposed that whatever positive value is attributed to either of these, the corresponding negative one is also given to it. $i = \sqrt{-1}$, $\omega$ and $\varpi$ are real positive quantities. (At least this is the standard case, and the only one we shall explicitly consider. Many of the formulae obtained are true, with slight modifications, whatever $\omega$ and $\varpi$ represent, provided only $\omega : \varpi$ be not a real quantity; for if it were so, $m\omega + n\varpi i$ for some values of $m$, $n$ would vanish, or at least become indefinitely small, and $u$ would cease to be a determinate function of $x$.)\footnote{I have examined the case of impossible values of $\omega$ and $\varpi$ in a paper which I am preparing for Crelle's Journal. [The paper here referred to is [25], actually published in Liouville's Journal].}

Now the value of the above expression, or, as for the sake of shortness it may be written, of the function
\[
u = x\prod\prod\left(1 + \frac{x}{\mn}\right) \cdots(3),
\]

% ===========================================================================
% Page 138
% ===========================================================================
\newpage
\setcounter{page}{138}

depends in a remarkable manner on the mode in which the superior limits of $m$, $n$ are assigned. Imagine $m$, $n$ to have any positive or negative integer values satisfying the equation
\[
\phi(m^2, n^2) < T \cdots(4).
\]
Consider, for greater distinctness, $m$, $n$ as the coordinates of a point; the equation $\phi(m^2, n^2) = T$ belongs to a certain curve symmetrical with respect to the two axes. I suppose besides that this is a continuous curve without multiple points, and such that the minimum value of a radius vector through the origin continually increases as $T$ increases, and becomes infinite with $T$. The curve may be analytically discontinuous, this is of no importance. The condition with respect to the limits is then that $m$ and $n$ must be integer values denoting the coordinates of a point \textit{within} the above curve, the whole system of such integer values being successively taken for these quantities.

Suppose, next, $u'$ denotes the same function as $u$, except that the limiting condition is
\[
\phi'(m^2, n^2) < T' \cdots(5).
\]
The curve $\phi'(m^2, n^2) = T'$ is supposed to possess the same properties with the other limiting curve, and, for greater distinctness, to lie entirely outside of it; but this last condition is non-essential.

These conditions being satisfied, the ratio $u' : u$ is very easily determined in the limiting case of $T$ and $T'$ infinite. In fact
\[
\frac{u'}{u} = \prod\prod\left(1 + \frac{x}{\mn}\right) \cdots(6),
\]
or
\[
l\frac{u'}{u} = \Sigma\Sigma\, l\!\left\{1 + \frac{x}{\mn}\right\} \cdots(7),
\]
the limiting conditions being
\[
\phi(m^2, n^2) > T,\quad \phi'(m^2, n^2) < T' \cdots(8).
\]

Now
\[
l\!\left\{1 + \frac{x}{\mn}\right\} = \frac{x}{\mn} - \tfrac{1}{2}\frac{x^2}{(\mn)^2} + \cdots(9),
\]
\[
l\frac{u'}{u} = x\cdot\Sigma\Sigma\frac{1}{\mn} - \tfrac{1}{2}x^2\cdot\Sigma\Sigma\frac{1}{(\mn)^2} + \cdots(10),
\]
or, the alternate terms vanishing on account of the positive and negative values destroying each other,
\[
l\frac{u'}{u} = -\tfrac{1}{2}x^2\cdot\Sigma\Sigma\frac{1}{(\mn)^2} - \tfrac{1}{4}x^4\cdot\Sigma\Sigma\frac{1}{(\mn)^4} - \cdots(11).
\]
In general
\[
\Sigma\Sigma\,\psi(m, n) = \iint\psi(m, n)\,dm\,dn + P \cdots(12),
\]

% ===========================================================================
% Page 139
% ===========================================================================
\newpage
\setcounter{page}{139}

$P$ denoting a series the first term of which is of the form $C\psi(m, n)$, and the remaining ones depending on the differential coefficients of this quantity with respect to $m$ and $n$. The limits between which the two sides are to be taken, are identical.

In the present case, supposing $T$ and $T'$ indefinitely great, it is easy to see that the first term of the expression for $l\frac{u'}{u}$ is the only one which is not indefinitely small; and we have
\[
l\frac{u'}{u} = -\tfrac{1}{2}Ax^2,\quad \text{or}\quad u' = u e^{-\frac{1}{2}Ax^2} \cdots(13),
\]
where
\[
A = \iint\frac{dm\,dn}{(\mn)^2} = \iint\frac{dm\,dn}{(m\omega + n\varpi i)^2} \cdots(14),
\]
the limits of the integration being given by
\[
\phi(m^2, n^2) > T,\quad \phi'(m^2, n^2) < T' \cdots(15).
\]
Some particular cases are important. Suppose the limits of $u'$ are given by
\[
m^2\omega^2 < T^2,\quad n^2\varpi^2 < T'^2 \cdots(16),
\]
and those of $u$, by
\[
m^2\omega^2 + n^2\varpi^2 < T^2 \cdots(17);
\]
we have
\[
A = \iint\frac{dm\,dn}{(m\omega + n\varpi i)^2}
\]
\[
= \frac{1}{\omega}\int_{-T'}^{T'}\!d(n\varpi)\!\left\{\frac{1}{T + n\varpi i} - \frac{1}{\sqrt{T^2 - n^2\varpi^2}+n\varpi i} - \frac{1}{-\sqrt{T^2 - n^2\varpi^2}+n\varpi i} - \frac{1}{-T+n\varpi i}\right\}
\]
\[
= -\frac{2}{\omega}\int_0^1 d\beta\!\left\{\sqrt{1-\beta^2}\right\} - \tfrac{2}{\omega}\,(\pi - \pi) = 0;
\]
or, in this case,
\[
u' = u \cdots(19).
\]
Again, let the limits of $u'$ be
\[
m^2\omega^2 < R'^2,\quad n^2\varpi^2 < S'^2 \cdots(20),
\]
and those of $u$,
\[
m^2\omega^2 < R^2,\quad n^2\varpi^2 < S^2 \cdots(21);
\]
\[
A = \iint\frac{dm\,dn}{(m\omega + n\varpi i)^2}
\]
\[
= \frac{1}{\omega}\int d(n\varpi)\!\left\{\frac{1}{R' + n\varpi i} - \frac{1}{R + n\varpi i} - \frac{1}{-R + n\varpi i} + \frac{1}{-R' + n\varpi i}\right\},
\]

% ===========================================================================
% Page 140
% ===========================================================================
\newpage
\setcounter{page}{140}

where the limits are $n^2\varpi^2 < S'^2$, for the terms containing $R'$, $n^2\varpi^2 < S^2$, for the terms containing $R$,
\[
= \frac{2}{\omega\varpi}\,l\!\left(\frac{R' + S'i}{R' - S'i}\cdot\frac{R - Si}{R + Si}\right) \cdots(23),
\]
\[
= -\frac{4}{\omega\varpi}\,i(\lambda' - \lambda),\quad \text{if}\ \lambda' = \tan^{-1}\frac{S'}{R'},\ \lambda = \tan^{-1}\frac{S}{R},
\]
the arcs $\lambda$, $\lambda'$ being included between the limits $0$, $\tfrac{1}{2}\pi$. Hence
\[
u' = u\,e^{2(\lambda' - \lambda)\beta x^2} \cdots(24).
\]
In particular if $\phi = \phi'$, $u' = u$. If $\frac{S}{R} = 0$, $\lambda = 0$, $u' = u\,e^{-i\beta x^2}$; if $\frac{S'}{R'} = \infty$, $\lambda' = \tfrac{1}{2}\pi$, $u' = u\,e^{i\beta x^2}$: where $\beta = \dfrac{\pi}{\omega\varpi}$, for which quantity it will continue to be used.

We may now completely define the functions whose properties are to be investigated. Writing, for shortness,
\begin{align*}
\mn &= m\omega + n\varpi i, \cdots(A)\\
\mbn &= (m + \tfrac{1}{2})\omega + n\varpi i,\\
\mnb &= m\omega + (n + \tfrac{1}{2})\varpi i,\\
\mbnb &= (m + \tfrac{1}{2})\omega + (n + \tfrac{1}{2})\varpi i;
\end{align*}
we may put
\begin{align*}
\gamma x &= x\prod\prod\!\left\{1 + \frac{x}{\mn}\right\}, \cdots(B)\\
g x &= \prod\prod\!\left\{1 + \frac{x}{\mbn}\right\},\\
G x &= \prod\prod\!\left\{1 + \frac{x}{\mnb}\right\},\\
\fS x &= \prod\prod\!\left\{1 + \frac{x}{\mbnb}\right\},
\end{align*}
the limits being given respectively by the equations
\[
\mathrm{mod.}\,\mn < T,\ \ \mathrm{mod.}\,\mbn < T,\ \ \mathrm{mod.}\,\mnb < T,\ \ \mathrm{mod.}\,\mbnb < T,
\]
$T$ being finally infinite. The system of values $m = 0$, $n = 0$, is of course omitted in $\gamma x$.

The functions $\gamma x$, $g x$, $G x$, $\fS x$, are all of them real finite functions of $x$, possessing properties analogous to that of $u$. Thus, representing any one of them by $Jx$, we have
\[
J x = e^{\frac{1}{2}i\beta\sigma x^2}\, J_{\pm\rho}\,x \cdots(C),
\]

% ===========================================================================
% Page 141
% ===========================================================================
\newpage
\setcounter{page}{141}

where $J_{\pm\rho}\,x$ is the same as $J x$, only for $J_\rho\,x$ the limits are given by $m^2\omega^2$ or $(m + \tfrac{1}{2})^2\omega^2 < R$, $n^2\varpi^2$ or $(n + \tfrac{1}{2})^2\varpi^2 < S$, ($R$, $S$, and $\frac{S}{R}$ infinite), and for $J_{-\rho}\,x$, by the same formulae, ($R$, $S$, and $\frac{R}{S}$ infinite). It is to this equation that the most characteristic properties of the functions $J x$ are due.

The following equations are deduced immediately from the above definitions:
\begin{align*}
\gamma(-x) &= -\gamma x, & g(-x) &= g x, & G(-x) &= G x, & \fS(-x) &= \fS x, \cdots(D)\\
\gamma(0) &= 0, & g(0) &= 1, & G(0) &= 1, & \fS(0) &= 1,\\
\gamma'(0) &= 1.
\end{align*}
Suppose $\gamma_1 x$, $g_1 x$, $G_1 x$, $\fS_1 x$, are the values that would have been obtained for $\gamma x$, $g x$, $G x$, $\fS x$ by interchanging $\omega$ and $\varpi$,---then changing $x$ into $x i$, and interchanging $m$ and $n$, by which means the limiting equations are the same in the two cases, we obtain the following system of equations:
\begin{align*}
\gamma_1(x i) &= i\,\gamma x, \cdots(E)\\
g_1(x i) &= G x,\\
G_1(x i) &= g x,\\
\fS_1(x i) &= \fS x;
\end{align*}
or otherwise,
\begin{align*}
\gamma(x i) &= i\,\gamma_1 x, \cdots(F)\\
g(x i) &= G_1 x,\\
G(x i) &= g_1 x,\\
\fS(x i) &= \fS_1 x,
\end{align*}
equations which are useful in transforming almost any other property of the functions $J$.

The functions $J x$ are changed one into another, except as regards a constant multiplier, by the change of $x$ into $x + \tfrac{1}{2}\omega$. This will be shown in a Note, or it may be seen from some formulae deduced immediately from the definitions of the functions $J$, which will be given in the sequel\footnote{Not given in the present paper. [The Note was given, see p.~154, and the formulas referred to must have been the formulae (M) p.~144.]}. Observing the relation between $J x$ and $J_\rho x$, we have in particular
\begin{align*}
\gamma\!\left(x + \tfrac{\omega}{2}\right) &= e^{\frac{1}{2}i\beta\omega x}\,A\,g x, \cdots(G)\\
g\!\left(x + \tfrac{\omega}{2}\right) &= e^{\frac{1}{2}i\beta\omega x}\,B\,\gamma x,\\
G\!\left(x + \tfrac{\omega}{2}\right) &= e^{\frac{1}{2}i\beta\omega x}\,C\,\fS x,\\
\fS\!\left(x + \tfrac{\omega}{2}\right) &= e^{\frac{1}{2}i\beta\omega x}\,D\,G x,
\end{align*}

% ===========================================================================
% Page 142
% ===========================================================================
\newpage
\setcounter{page}{142}

where $A$, $B$, $C$, $D$ are most simply determined by writing $x = 0$, $x = -\tfrac{\omega}{2}$. Putting at the same time $e^{i\beta\omega^2/4} = e^{\pi\omega/4\varpi} = q^{-\frac{1}{4}}$,
\begin{align*}
A &= \gamma\!\left(\tfrac{\omega}{2}\right), \cdots(H)\\
B &= -q^{-1} \div \gamma\!\left(\tfrac{\omega}{2}\right),\\
C &= G\!\left(\tfrac{\omega}{2}\right) = q^{-1}\cdot\fS^{-1}\!\!\left(\tfrac{\omega}{2}\right),\\
D &= \fS\!\left(\tfrac{\omega}{2}\right) = q^{-1}\cdot G^{-1}\!\!\left(\tfrac{\omega}{2}\right);
\end{align*}
whence also
\[
G\!\left(\tfrac{\omega}{2}\right)\fS\!\left(\tfrac{\omega}{2}\right) = q^{-1} \cdots(25).
\]
Similarly, the functions $J x$ are changed one into the other by the change of $x$ into $x + \tfrac{1}{2}\varpi i$. We have in the same way
\begin{align*}
\gamma\!\left(x + \tfrac{\varpi i}{2}\right) &= e^{-\frac{1}{2}i\beta\varpi x}\,A'\,G x, \cdots(I)\\
g\!\left(x + \tfrac{\varpi i}{2}\right) &= e^{-\frac{1}{2}i\beta\varpi x}\,B'\,\fS x,\\
G\!\left(x + \tfrac{\varpi i}{2}\right) &= e^{-\frac{1}{2}i\beta\varpi x}\,C'\,\gamma x,\\
\fS\!\left(x + \tfrac{\varpi i}{2}\right) &= e^{-\frac{1}{2}i\beta\varpi x}\,D'\,g x.
\end{align*}
Whence
\begin{align*}
A' &= \gamma\!\left(\tfrac{\varpi i}{2}\right), \cdots(J)\\
B' &= g\!\left(\tfrac{\varpi i}{2}\right) = q_1^{-1} \div \fS\!\left(\tfrac{\varpi i}{2}\right),\\
C' &= -q_1^{-1} \div \gamma\!\left(\tfrac{\varpi i}{2}\right),\\
D' &= \fS\!\left(\tfrac{\varpi i}{2}\right) = q_1^{-1} \div g\!\left(\tfrac{\varpi i}{2}\right);
\end{align*}
where $e^{i\beta\varpi^2/4} = e^{\pi\varpi/4\omega} = q_1^{-\frac{1}{4}}$. It is obvious that the relation between $q$ and $q_1$ is $lq\cdot lq_1 = -\pi^2$.

We obtain from the above
\[
g\!\left(\tfrac{\varpi i}{2}\right)\fS\!\left(\tfrac{\varpi i}{2}\right) = q_1^{-1} \cdots(26).
\]

% ===========================================================================
% Page 143
% ===========================================================================
\newpage
\setcounter{page}{143}

Also, by making $x = \tfrac{\varpi i}{2}$ in the expression for $\gamma\!\left(x + \tfrac{\omega}{2}\right)$ and $x = \tfrac{\omega}{2}$ in that for $\gamma\!\left(x + \tfrac{\varpi i}{2}\right)$, we have
\[
\gamma\!\left(\tfrac{\omega}{2}\right)g\!\left(\tfrac{\varpi i}{2}\right) = -i\,\gamma\!\left(\tfrac{\varpi i}{2}\right)G\!\left(\tfrac{\omega}{2}\right) \cdots(27),
\]
and the same or an equivalent one would have been obtained from the functions $g$, $G$, $\fS$.

By combining the above systems, we deduce one of the form
\begin{align*}
\gamma\!\left(x + \tfrac{\omega}{2} + \tfrac{\varpi i}{2}\right) &= e^{i\beta x(\omega - \varpi i)}\,A''\,\fS x, \cdots(K)\\
g\!\left(x + \tfrac{\omega}{2} + \tfrac{\varpi i}{2}\right) &= e^{i\beta x(\omega - \varpi i)}\,B''\,G x,\\
G\!\left(x + \tfrac{\omega}{2} + \tfrac{\varpi i}{2}\right) &= e^{i\beta x(\omega - \varpi i)}\,C''\,g x,\\
\fS\!\left(x + \tfrac{\omega}{2} + \tfrac{\varpi i}{2}\right) &= e^{i\beta x(\omega - \varpi i)}\,D''\,\gamma x,
\end{align*}
and, observing the equation $e^{i\beta\omega\varpi i/4} = e^{\pi/4} = (-1)^{1/4}$, with the following values for the coefficients,
\begin{align*}
A'' &= -(-1)^{\frac{1}{4}}\cdot\gamma\!\left(\tfrac{\omega}{2}\right)\cdot\fS\!\left(\tfrac{\omega}{2}\right), \cdots(L)\\
B'' &= -(-1)^{\frac{1}{4}}\cdot q^{-1}\cdot\gamma\!\left(\tfrac{\omega}{2}\right)\div g\!\left(\tfrac{\omega}{2}\right),\\
C'' &= -(-1)^{\frac{1}{4}}\cdot\fS\!\left(\tfrac{\omega}{2}\right)\cdot G\!\left(\tfrac{\omega}{2}\right),\\
D'' &= -(-1)^{\frac{1}{4}}\cdot G\!\left(\tfrac{\omega}{2}\right)\cdot\fS\!\left(\tfrac{\omega}{2}\right).
\end{align*}

Collecting the formulae which connect $\gamma\!\left(\tfrac{\omega}{2}\right)$, $\gamma\!\left(\tfrac{\varpi i}{2}\right)$\ldots\ these are
\begin{align*}
\gamma\!\left(\tfrac{\omega}{2}\right) &= 0, \cdots(L\,bis)\\
G\!\left(\tfrac{\omega}{2}\right)\fS\!\left(\tfrac{\omega}{2}\right) &= q^{-1},\\
g\!\left(\tfrac{\varpi i}{2}\right)\fS\!\left(\tfrac{\varpi i}{2}\right) &= q_1^{-1},\\
\gamma\!\left(\tfrac{\omega}{2}\right)g\!\left(\tfrac{\varpi i}{2}\right) &= -i\,\gamma\!\left(\tfrac{\varpi i}{2}\right)G\!\left(\tfrac{\omega}{2}\right).
\end{align*}

% ===========================================================================
% Page 144
% ===========================================================================
\newpage
\setcounter{page}{144}

And by the assistance of these
\begin{align*}
B'' C'' \div A'' D'' &= B' D' \div A' C' = C D \div A B = -1 \cdots(28),\\
A' B' \div C' D' &= -A'' B'' \div C'' D'' = -\gamma^2\!\left(\tfrac{\omega}{2}\right) \div \fS^2\!\left(\tfrac{\omega}{2}\right),\\
A'' C'' \div B'' D'' &= -A' C' \div B' D' = -\gamma^2\!\left(\tfrac{\varpi i}{2}\right) \div \fS^2\!\left(\tfrac{\varpi i}{2}\right),\\
A D \div B C &= -A' D' \div B' C' = -\gamma\!\left(\tfrac{\omega}{2}\right) \div \fS\!\left(\tfrac{\omega}{2}\right) = -\gamma\!\left(\tfrac{\varpi i}{2}\right) \div \fS\!\left(\tfrac{\varpi i}{2}\right),
\end{align*}
which will be required presently.

It is now easy to proceed to the general systems of formulae,
\begin{align*}
\Theta &= (-1)^{mn}\,e^{i\beta x(m\omega - n\varpi i)}\,q_1^{-\frac{1}{4}m^2}\,q^{-\frac{1}{4}n^2}, \cdots(M)\\
\gamma\{x + \mn\} &= (-1)^{m+n}\cdot\Theta\,\gamma x,\\
g\{x + \mn\} &= (-1)^m \cdot\Theta\,g x,\\
G\{x + \mn\} &= (-1)^n \cdot\Theta\,G x,\\
\fS\{x + \mn\} &= \Theta\,\fS x.
\end{align*}
\begin{align*}
\Phi &= (-1)^{m+n}\,e^{i\beta x[(m+\frac{1}{2})\omega - n\varpi i]}\,q_1^{-\frac{1}{4}(m+\frac{1}{2})^2}\,q^{-\frac{1}{4}n^2}, \cdots\\
\gamma\{x + \mbn\} &= (-1)^{m+n}\cdot\Phi\,A\,g x,\\
g\{x + \mbn\} &= (-1)^m \cdot\Phi\,B\,\gamma x,\\
G\{x + \mbn\} &= (-1)^n \cdot\Phi\,C\,\fS x,\\
\fS\{x + \mbn\} &= \Phi\,D\,G x.
\end{align*}
\begin{align*}
\Psi &= (-1)^{m+n}\,e^{i\beta x[m\omega - (n+\frac{1}{2})\varpi i]}\,q_1^{-\frac{1}{4}m^2}\,q^{-\frac{1}{4}(n+\frac{1}{2})^2}, \cdots\\
\gamma\{x + \mnb\} &= (-1)^{m+n}\cdot\Psi\,A'\,G x,\\
g\{x + \mnb\} &= (-1)^m \cdot\Psi\,B'\,\fS x,\\
G\{x + \mnb\} &= (-1)^n \cdot\Psi\,C'\,\gamma x,\\
\fS\{x + \mnb\} &= \Psi\,D'\,g x.
\end{align*}
\begin{align*}
\Omega &= (-1)^{m+n+1}\,e^{i\beta x[(m+\frac{1}{2})\omega - (n+\frac{1}{2})\varpi i]}\,q_1^{-\frac{1}{4}(m+\frac{1}{2})^2}\,q^{-\frac{1}{4}(n+\frac{1}{2})^2}, \cdots\\
\gamma\{x + \mbnb\} &= (-1)^{m+n+1}\cdot\Omega\,A''\,\fS x,\\
g\{x + \mbnb\} &= (-1)^m \cdot\Omega\,B''\,G x,\\
G\{x + \mbnb\} &= (-1)^n \cdot\Omega\,C''\,g x,\\
\fS\{x + \mbnb\} &= \Omega\,D''\,\gamma x.
\end{align*}

% ===========================================================================
% Page 145
% ===========================================================================
\newpage
\setcounter{page}{145}

Suppose $x = 0$, we have the new systems,
\begin{align*}
\Theta_0 &= (-1)^{mn}\,q_1^{-\frac{1}{4}m^2}\,q^{-\frac{1}{4}n^2}, \cdots(M\,bis)\\
\gamma\mn &= 0,\quad &\gamma'\mn &= (-1)^{m+n}\,\Theta_0,\\
g\mn &= (-1)^m\,\Theta_0,\\
G\mn &= (-1)^n\,\Theta_0,\\
\fS\mn &= \Theta_0,
\end{align*}
\begin{align*}
\Phi_0 &= (-1)^{m+n}\,q_1^{-\frac{1}{4}(m+\frac{1}{2})^2}\,q^{-\frac{1}{4}n^2}, \cdots\\
\gamma\mbn &= (-1)^{m+n}\,\Phi_0\,A,\\
g\mbn &= 0,\quad &g'\mbn &= (-1)^m\,\Phi_0\,B,\\
G\mbn &= (-1)^n\,\Phi_0\,C,\\
\fS\mbn &= \Phi_0\,D.
\end{align*}
\begin{align*}
\Psi_0 &= (-1)^{m+n}\,q_1^{-\frac{1}{4}m^2}\,q^{-\frac{1}{4}(n+\frac{1}{2})^2}, \cdots\\
\gamma\mnb &= (-1)^{m+n}\,\Psi_0\,A',\\
g\mnb &= (-1)^m\,\Psi_0\,B',\\
G\mnb &= 0,\quad &G'\mnb &= (-1)^n\,\Psi_0\,C',\\
\fS\mnb &= \Psi_0\,D'.
\end{align*}
\begin{align*}
\Omega_0 &= (-1)^{m+n+1}\,q_1^{-\frac{1}{4}(m+\frac{1}{2})^2}\,q^{-\frac{1}{4}(n+\frac{1}{2})^2}, \cdots\\
\gamma\mbnb &= (-1)^{m+n+1}\,\Omega_0\,A'',\\
g\mbnb &= (-1)^m\,\Omega_0\,B'',\\
G\mbnb &= (-1)^n\,\Omega_0\,C'',\\
\fS\mbnb &= 0,\quad &\fS'\mbnb &= \Omega_0\,D''.
\end{align*}

We obtain immediately, by taking the logarithmic differentials of the functions $\gamma x$, $g x$, $G x$, $\fS x$, the equations
\begin{align*}
\gamma' x \div \gamma x &= \Sigma\Sigma\{x - \mn\}^{-1}, \quad m=0,\ n=0\ \text{admissible}, \cdots(N)\\
g' x \div g x &= \Sigma\Sigma\{x - \mbn\}^{-1},\\
G' x \div G x &= \Sigma\Sigma\{x - \mnb\}^{-1},\\
\fS' x \div \fS x &= \Sigma\Sigma\{x - \mbnb\}^{-1},
\end{align*}
the limits being the same as in the case of the factorial expressions.

Consider an equation
\[
g x\,G x \div \gamma x\,\fS x = \Sigma\Sigma\!\left[\fA\{x - \mn\}^{-1} + \fB\{x - \mbn\}^{-1}\right] \cdots(29),
\]

% ===========================================================================
% Page 146
% ===========================================================================
\newpage
\setcounter{page}{146}

we have
\[
\fA = g\mn\,G\mn \div \gamma'\mn\,\fS\mn = 1 \cdots(30),
\]
\[
\fB = g\mbn\,G\mbn \div \gamma\mbn\,\fS'\mbn = B'C' \div A''D'' = -1 \cdots(31).
\]
(The application of the ordinary method of decomposition into partial fractions, which is in general exceedingly precarious when applied to transcendental functions, is justified here by a theorem of Cauchy's, which will presently be quoted.) We have thus
\[
g x\,G x \div \gamma x\,\fS x = (\gamma' x \div \gamma x) - (\fS' x \div \fS x),
\]
and similarly
\begin{align*}
g x\,\fS x \div \gamma x\,G x &= (\gamma' x \div \gamma x) - (G' x \div G x), \cdots(O)\\
G x\,\fS x \div \gamma x\,g x &= (\gamma' x \div \gamma x) - (g' x \div g x),\\
-\gamma x\,\fS x \div g x\,G x &= (g' x \div g x) - (G' x \div G x),\\
\gamma x\,g x \div G x\,\fS x &= (G' x \div G x) - (\fS' x \div \fS x),\\
\gamma x\,G x \div g x\,\fS x &= (g' x \div g x) - (\fS' x \div \fS x);
\end{align*}
in which we have written
\[
\gamma\!\left(\tfrac{\omega}{2}\right)\!\div\!\fS\!\left(\tfrac{\omega}{2}\right) = c,\quad \gamma\!\left(\tfrac{\varpi i}{2}\right)\!\div\!\fS\!\left(\tfrac{\varpi i}{2}\right) = \tfrac{i}{e} \cdots(32).
\]
Eliminating the derived coefficients,
\begin{align*}
G^2 x - \fS^2 x &= c^2 \gamma^2 x, \cdots(33)\\
g^2 x - G^2 x &= -b^2 \gamma^2 x,\\
\fS^2 x - g^2 x &= \tfrac{1}{e^2}\,\gamma^2 x.
\end{align*}
Adding these equations, $b^2 = c^2 + \tfrac{1}{e^2}$, or $b = \sqrt{(c^2 + \tfrac{1}{e^2})}$, in which sense it will continue to be used.

Also,
\begin{align*}
g^2 x &= \fS^2 x - c^2 \gamma^2 x, \cdots(P)\\
G^2 x &= \fS^2 x + \tfrac{1}{e^2}\gamma^2 x.
\end{align*}
Suppose
\begin{align*}
\phi x &= \gamma x \div \fS x, \cdots(Q)\\
f x &= g x \div \fS x,\\
F x &= G x \div \fS x,
\end{align*}
then
\begin{align*}
f^2 x &= 1 - c^2 \phi^2 x, \cdots(R)\\
F^2 x &= 1 + \tfrac{1}{e^2}\phi^2 x,
\end{align*}

% ===========================================================================
% Page 147
% ===========================================================================
\newpage
\setcounter{page}{147}

and also
\begin{align*}
\phi' x &= f x \cdot F x, \cdots(S)\\
f' x &= -c^2 \phi x\,F x,\\
F' x &= \tfrac{1}{e^2}\phi x\,f x.
\end{align*}
Hence, putting for $f x$, $F x$, their values,
\[
1 = \frac{\phi' x}{\sqrt{(1 - c^2 \phi^2 x)(1 + \tfrac{1}{e^2}\phi^2 x)}} \cdots(T),
\]
or writing $\phi x = y$, and integrating,
\[
x = \int_0^y \frac{dy}{\sqrt{(1 - c^2 y^2)(1 + \tfrac{1}{e^2} y^2)}} \cdots(U),
\]
or
\[
\phi^{-1} y = \int_0^y \frac{dy}{\sqrt{(1 - c^2 y^2)(1 + \tfrac{1}{e^2} y^2)}},
\]
which shows that $\phi$ is an inverse elliptic function.

The equations which are the foundation of the theory of the functions $\phi$, $f$, $F$, are deduced immediately from the equations (S). (Abel, \OE uvres, tom.~i.\ p.~143 [Ed.~2, p.~268.]) These are
\begin{align*}
\phi(x+y) &= \frac{\phi x\,f y\,F y + \phi y\,f x\,F x}{1 + c^2 e^2\,\phi^2 x\,\phi^2 y}, \cdots(V)\\
f(x+y) &= \frac{f x\,f y - c^2\,\phi x\,\phi y\,F x\,F y}{1 + c^2 e^2\,\phi^2 x\,\phi^2 y},\\
F(x+y) &= \frac{F x\,F y + e^2\,\phi x\,\phi y\,f x\,f y}{1 + c^2 e^2\,\phi^2 x\,\phi^2 y};
\end{align*}
so that from this point we may take for granted any properties of these functions. We see, for instance, immediately,
\[
\phi\!\left(\tfrac{\omega}{2}\right) = \tfrac{1}{c},\quad \phi\!\left(\tfrac{\varpi i}{2}\right) = \tfrac{i}{e};
\]
whence
\[
\tfrac{\omega}{2} = \int_0^{\frac{1}{c}}\!\frac{dy}{\sqrt{(1 - c^2 y^2)(1 + \tfrac{1}{e^2} y^2)}} \cdots(W),
\]
\[
\tfrac{\varpi i}{2} = \int_0^{\frac{i}{e}}\!\frac{dy}{\sqrt{(1 - c^2 y^2)(1 + \tfrac{1}{e^2} y^2)}},\quad \text{or}\quad \tfrac{\varpi}{2} = \int_0^{\frac{1}{e}}\!\frac{dy}{\sqrt{(1 + c^2 y^2)(1 - \tfrac{1}{e^2} y^2)}} \cdots(X),
\]
which give the values of $\omega$, $\varpi$ in terms of $c$, $e$; values which may be developed in a variety of ways, in infinite series. We may also express $\gamma\!\left(\tfrac{\omega}{2}\right)$, \&c., and consequently $A$, $B$,\ldots\&c., by means of the quantities $c$, $e$. We have only to combine the equations
\[
\gamma\!\left(\tfrac{\omega}{2}\right)\div\fS\!\left(\tfrac{\omega}{2}\right) = c,\ \ G\!\left(\tfrac{\omega}{2}\right)\fS\!\left(\tfrac{\omega}{2}\right) = q^{-1},\ \ \fS\!\left(\tfrac{\varpi i}{2}\right)\div\gamma\!\left(\tfrac{\varpi i}{2}\right) = \tfrac{e}{i},\ \ g\!\left(\tfrac{\varpi i}{2}\right)\fS\!\left(\tfrac{\varpi i}{2}\right) = q_1^{-1},
\]

% ===========================================================================
% Page 148
% ===========================================================================
\newpage
\setcounter{page}{148}

with the former relations between these quantities, and we have
\begin{align*}
\gamma\!\left(\tfrac{\omega}{2}\right) &= b^{-\frac{1}{2}} c^{-\frac{1}{2}} q^{-\frac{1}{4}},\quad &\gamma\!\left(\tfrac{\varpi i}{2}\right) &= ib^{-\frac{1}{2}} e^{\frac{1}{2}} q_1^{-\frac{1}{4}}, \cdots(Y)\\
g\!\left(\tfrac{\omega}{2}\right) &= 0,\quad &g\!\left(\tfrac{\varpi i}{2}\right) &= b^{\frac{1}{2}} e^{-\frac{1}{2}} q_1^{-\frac{1}{4}},\\
G\!\left(\tfrac{\omega}{2}\right) &= b^{\frac{1}{2}} c^{-\frac{1}{2}} q^{-\frac{1}{4}},\quad &G\!\left(\tfrac{\varpi i}{2}\right) &= 0,\\
\fS\!\left(\tfrac{\omega}{2}\right) &= b^{-\frac{1}{2}} c^{-\frac{1}{2}} q^{-\frac{1}{4}},\quad &\fS\!\left(\tfrac{\varpi i}{2}\right) &= b^{-\frac{1}{2}} e^{\frac{1}{2}} q_1^{-\frac{1}{4}};
\end{align*}
\begin{align*}
A &= b^{-\frac{1}{2}} c^{-\frac{1}{2}} q^{-\frac{1}{4}}, \cdots(Z)\\
B &= -b^{\frac{1}{2}} c^{\frac{1}{2}} q^{\frac{1}{4}},\\
C &= b^{\frac{1}{2}} c^{-\frac{1}{2}} q^{-\frac{1}{4}},\\
D &= b^{-\frac{1}{2}} c^{\frac{1}{2}} q^{\frac{1}{4}},\\[4pt]
A' &= ib^{-\frac{1}{2}} e^{\frac{1}{2}} q_1^{-\frac{1}{4}},\quad &A'' &= (-1)^{\frac{1}{4}} c^{-\frac{1}{2}} e^{\frac{1}{2}} q^{-\frac{1}{4}} q_1^{-\frac{1}{4}},\\
B' &= b^{\frac{1}{2}} e^{-\frac{1}{2}} q_1^{-\frac{1}{4}},\quad &B'' &= -(-1)^{\frac{1}{4}} c^{\frac{1}{2}} e^{-\frac{1}{2}} q^{\frac{1}{4}} q_1^{-\frac{1}{4}},\\
C' &= -ib^{\frac{1}{2}} e^{-\frac{1}{2}} q_1^{\frac{1}{4}},\quad &C'' &= (-1)^{\frac{1}{4}} c^{-\frac{1}{2}} e^{\frac{1}{2}} q^{-\frac{1}{4}} q_1^{\frac{1}{4}},\\
D' &= b^{-\frac{1}{2}} e^{\frac{1}{2}} q_1^{\frac{1}{4}},\quad &D'' &= -(-1)^{\frac{1}{4}} c^{\frac{1}{2}} e^{-\frac{1}{2}} q^{\frac{1}{4}} q_1^{\frac{1}{4}},
\end{align*}
which are to be substituted in any formulae into which these quantities enter.

The following is Cauchy's Theorem, (\textit{Exercises de Math.\ t.\ ii.\ p.~289}).

``If in attributing to the modulus $r$ of the variable
\[
z = r(\cos p + \sqrt{(-1)}\sin p) \cdots(35),
\]
infinitely great values, these can be chosen so that the two functions
\[
\frac{f z + f(-z)}{2},\quad \frac{f z - f(-z)}{2 z} \cdots(36),
\]
sensibly vanish, whatever be the value of $p$, or vanish in general, though ceasing to do so and obtaining \textit{finite} values for certain particular values of $p$; then
\[
f x = \sum\frac{(f z)}{z - x} \cdots(37),
\]
the integral residue being reduced to its principal value.''

To understand this, it is only necessary to remark that the integral residue in question is the series of fractions that would be obtained by the ordinary process of decomposition; and by the principal value is meant, that all those roots are to be taken, the modulus of which is not greater than a certain limit, this limit being afterwards made infinite.

Suppose now $f x$ is a fraction, the numerator and denominator of which are monomials of the form $(\gamma x)^l (g x)^m \ldots$, $l, m, \ldots$ being positive integers, and of course no common factor being left in the numerator and denominator.

Let $\lambda$ be the excess of the degree of the denominator over that of the numerator. Suppose the modulus $r$ of $(z)$ has any value not the same with any of the moduli of
\[
\mn,\ \mbn,\ \mnb,\ \mbnb \cdots(38).
\]

% ===========================================================================
% Page 149
% ===========================================================================
\newpage
\setcounter{page}{149}

Then we have
\[
r(\cos p + i\sin p) = m\omega + n\varpi i + \theta \cdots(39),
\]
$\theta$ being a finite quantity, such that none of the functions $J\theta$ vanish, $m$ and $n$ are the greatest integer values which allow the possible part of $\theta$ and the coefficient of its impossible part to remain positive. We have therefore
\[
m^2\omega^2 + n^2\varpi^2 = r^2 - M \cdots(40),
\]
$M$ being finite; or when $r$ is infinite, at least one of the values $m$, $n$ is infinite. The function $f z$ reduces itself to the form
\[
f z = F\,q_1^{\frac{1}{4}\lambda m^2}\,q^{\frac{1}{4}\lambda n^2}\,e^{\frac{1}{2}\lambda(m\omega + n\varpi i)} \cdots(41),
\]
where $F$ is finite. Hence $q$, and $q_1$ being always less than unity, $f z$, and consequently both $\tfrac{1}{2}\{f z + f(-z)\}$ and $\tfrac{1}{2}\{f z - f(-z)\}$ vanish for $r = \infty$, as long as $\lambda$ is positive.

In the case of $\lambda = 0$, the conditions are still satisfied, if we suppose $f x$ to denote an uneven function of $x$: for when $\lambda = 0$, the index of exponential in the above expression vanishes, or $f z$ is constantly finite. But $f z$ being an odd function of $z$, $f z + f(-z) = 0$. And $\tfrac{1}{z}[f z - f(-z)]$ vanishes for $z$ infinite, on account of the $z$ in the denominator: hence the expansion is admissible in this case. But it is certainly so also, in a great many cases at least, where $f z$ is an even function of $z$; for these may be deduced from the others by a simple change in the value of the variable. For instance, from the expansion of $\gamma x \div g x$, which is an odd function, by writing $x + \tfrac{\omega}{2}$ for $x$, we obtain that of $G x \div \fS x$, which is even.

A case of some importance is when the function is of the above form, multiplied by an exponential $e^{i a x^2 + b x}$. Here writing $z = m\omega + n\varpi i + \theta$, the admissibility of the formula depends on the evanescence of
\[
e^{i a(m\omega + n\varpi i)^2}\,q_1^{\frac{1}{4}\lambda m^2}\,q^{\frac{1}{4}\lambda n^2} \cdots(42);
\]
or, if $a = h + ki$, this becomes, omitting a finite factor,
\[
e^{-i m^2 i(\lambda\beta - h)\omega^2 - i n^2 i(\lambda\beta + h)\varpi^2 - 2 k m n\,\omega\varpi} \cdots(43),
\]
which vanishes if $h^2 + k^2 < \lambda^2\beta^2$, i.e.\ the modulus of $a$ is less than $\lambda\beta$. The limiting case is admissible when the series is convergent.

We obtain in this way a very great variety of formulae. For instance,
\begin{align*}
e^{iax^2 + bx}\div\gamma x &= \Sigma\Sigma\!\left[(-1)^{-mn-m-n}\,e^{i a(m\omega - n\varpi i)^2 + b(m\omega, n\varpi i)}\,q^{\frac{1}{4} a m^2}\,q^{\frac{1}{4} a n^2}\!\{x - \mn\}^{-1}\right] \cdots(A')\\
e^{iax^2 + bx}\div g x &= -b^{-\frac{1}{2}} c^{-\frac{1}{2}} \Sigma\Sigma\!\left[(-1)^{-\frac{1}{2}m+\frac{1}{2}-m-n}\,e^{i a(m+\frac{1}{2}\omega - n\varpi i)^2 + b(\ldots)}\,q_1^{\frac{1}{4} a (m+\frac{1}{2})^2}\,q^{\frac{1}{4} a n^2}\!\{x - \mbn\}^{-1}\right],\\
e^{iax^2 + bx}\div G x &= ib^{-\frac{1}{2}} e^{-\frac{1}{2}} \Sigma\Sigma\!\left[(-1)^{-m+\frac{1}{2}n-m-n}\,e^{i a(m\omega - (n+\frac{1}{2})\varpi i)^2 + b(\ldots)}\,q_1^{\frac{1}{4} a m^2}\,q^{\frac{1}{4} a (n+\frac{1}{2})^2}\!\{x - \mnb\}^{-1}\right],\\
e^{iax^2 + bx}\div \fS x &= -i c^{-\frac{1}{2}} e^{-\frac{1}{2}} \Sigma\Sigma\!\left[(-1)^{-\frac{1}{2}(m+1)(n+\frac{1}{2})}\,e^{i a(\ldots)+b(\ldots)}\,q_1^{\frac{1}{4} a (m+\frac{1}{2})^2}\,q^{\frac{1}{4} a (n+\frac{1}{2})^2}\!\{x - \mbnb\}^{-1}\right],
\end{align*}

% ===========================================================================
% Page 150
% ===========================================================================
\newpage
\setcounter{page}{150}

in which the modulus of $a$ must not exceed $\beta$: in the limiting cases, for $a = \beta$, $b$ must be entirely impossible, and for $a = -\beta$, $b$ must be entirely real. The formulae for $\gamma x$ are
\begin{align*}
e^{i\beta x^2 + b x}\div \gamma x &= \Sigma\Sigma\,(-1)^{-mn}\,q_1^{-\frac{1}{2}m^2}\,e^{b\mn}\{x - \mn\}^{-1} \cdots(44),\\
e^{-i\beta x^2 + b x}\div \gamma x &= \Sigma\Sigma\,(-1)^{-m-n}\,q^{-\frac{1}{2}n^2}\,e^{b\mn}\{x - \mn\}^{-1};
\end{align*}
and for $b = 0$,
\begin{align*}
e^{i\beta x^2}\div \gamma x &= \Sigma\Sigma\,(-1)^{-mn}\,q_1^{-\frac{1}{2}m^2}\{x - \mn\}^{-1} \cdots(45),\\
e^{-i\beta x^2}\div \gamma x &= \Sigma\Sigma\,(-1)^{-m-n}\,q^{-\frac{1}{2}n^2}\{x - \mn\}^{-1}.
\end{align*}
Next the system,
\begin{align*}
\fS x \div \gamma x &= \Sigma\Sigma\,(-1)^{m+n}\,\{x - \mn\}^{-1}, \cdots(B')\\
G x \div \gamma x &= \Sigma\Sigma\,(-1)^m\,\{x - \mn\}^{-1},\\
g x \div \gamma x &= \Sigma\Sigma\,(-1)^n\,\{x - \mn\}^{-1};\\[4pt]
\gamma x \div g x &= -b^{-1} c^{-1}\,\Sigma\Sigma\,(-1)^n\,\{x - \mbn\}^{-1},\\
G x \div g x &= -c^{-1}\,\Sigma\Sigma\,(-1)^{m+n}\,\{x - \mbn\}^{-1},\\
\fS x \div g x &= -b^{-1}\,\Sigma\Sigma\,(-1)^m\,\{x - \mbn\}^{-1};\\[4pt]
\gamma x \div G x &= -b^{-1} e^{-1}\,\Sigma\Sigma\,(-1)^m\,\{x - \mnb\}^{-1},\\
g x \div G x &= i e^{-1}\,\Sigma\Sigma\,(-1)^{m+n}\,\{x - \mnb\}^{-1},\\
\fS x \div G x &= i e^{-1}\,\Sigma\Sigma\,(-1)^n\,\{x - \mnb\}^{-1};\\[4pt]
\gamma x \div \fS x &= ic^{-1} e^{-1}\,\Sigma\Sigma\,(-1)^{m+n}\,\{x - \mbnb\}^{-1},\\
g x \div \fS x &= e^{-1}\,\Sigma\Sigma\,(-1)^m\,\{x - \mbnb\}^{-1},\\
G x \div \fS x &= i c^{-1}\,\Sigma\Sigma\,(-1)^n\,\{x - \mbnb\}^{-1};
\end{align*}
which is partially given by Abel.

We may obtain, in like manner, expressions for the functions
\[
\frac{1}{\gamma x\,g x},\ \frac{1}{\gamma x\,G x},\ldots\ (\text{six terms of this form}) \cdots(C'),
\]
\[
\frac{1}{\gamma x\,g x\,G x},\ldots\ (\text{four terms}) \cdots(D'),
\]
\[
\frac{1}{\gamma x\,g x\,G x\,\fS x} \cdots(E'),
\]
\[
\frac{1}{\gamma x\,g x\,G x\,\fS x^2} \cdots(F'),
\]
and similar systems with higher powers in the denominator.

\end{document}
